\section{Where the Stiffness Might Come From}

I left a debt standing a few pages ago and I would rather report on it than let it sit. The ladder's whole force depends on a strut I could not account for: degeneracy pressure, which the standard picture explains by a prohibition, and which this book may not explain that way, because a prohibition with nothing doing the forbidding is the same species of object as a pull across empty space. So I went looking at the one place in nature where circulations are packed together and \emph{measured}. What is there does not settle the question. It does narrow it, and it narrows it to a single line, which is worth more than a page of hedging.

Cool helium far enough and it becomes a superfluid, and it acquires a strange refusal: set the vessel spinning and the liquid declines to spin with it. What it does instead is thread itself with whirlpools --- vortices, each a line with the fluid turning around it --- and the whirlpools take up the rotation between them. The circulation of each one, meaning the flow summed around any loop enclosing it, is not roughly the same from vortex to vortex. It is \emph{exactly} the same: a fixed quantum, $\kappa = h/m$, about a ten-millionth of a square metre per second in helium, with $h$ Planck's constant and $m$ the mass of one atom. You cannot have two thirds of a vortex. Around each, the fluid turns at $\kappa/2\pi r$, and at the centre is a core about an \aa ngstr\"om across where the fluid simply gives out.\footnote{The core radius is set by the \emph{healing length}: the distance over which a superfluid can change its density. Inside it the superfluid has been destroyed; the vortex is a hole with a whirl around it.}

Now put that quantum to work, because it hands over the very thing the strut needs. Confine a circulation to a region of size $d$ and its speed there must be $\kappa/2\pi d$; so its momentum is
\[
p \;=\; m v \;=\; \frac{m\kappa}{2\pi d} \;=\; \frac{h}{2\pi d} \;=\; \frac{\hbar}{d},
\]
and notice what fell out of the arithmetic on the way: the mass cancelled. Squeeze a circulation into half the room and it must turn twice as fast, for the plainest reason there is --- the circulation it carries is fixed, and a fixed circulation in a smaller loop is a higher speed. That is precisely the confinement relation that degeneracy pressure runs on, and here it arrives with a mechanism under it instead of a decree.

Follow it one step further. Pack such circulations one to a cell of size $d$, so that the number in a unit of volume is $n = 1/d^{3}$. Each carries momentum $\hbar/d$, hence an energy of order $\hbar^{2}n^{2/3}/2m$; multiply by the number present and the energy in a unit of volume grows as $n^{5/3}$. That is the degenerate exponent \emph{exactly} --- the five-thirds that holds up white dwarfs --- and with the crudest geometry anyone could write down the size of it lands within a factor of about six of the real thing.

I want to say what that does and does not prove, immediately, before the pleasure of it does any damage. It does not prove exclusion. The calculation assumed one circulation to a cell, and that assumption \emph{is} the exclusion; nothing came out that was not put in. What has changed is which half of the problem is missing. The kinematics --- why crowding forces motion, and forces it in exactly the proportion that gives the observed stiffness --- is now free, and comes from a fixed quantum of circulation rather than from a rule. The packing is not free. Why one to a cell?

The fluid has something to say about that too, and it is more than I expected. Start from two facts that are not in dispute: the energy of a circulation goes as the \emph{square} of its strength, and circulations add when seen from far enough away. Put two of the same sense side by side and, from a distance, they look like a single circulation of twice the strength: four units of energy, where the two of them apart cost only two. So a double quantum is not forbidden --- it is simply a bad bargain, and it splits. Which is exactly what doubly-quantized vortices are observed to do. That is not a prohibition against two things sharing a place. It is an arithmetic that drives them apart, which is what a mechanical account of exclusion would have to look like.

And the two senses are not symmetric. Opposite circulations \emph{cancel} at a distance rather than adding, so a pair of them costs almost nothing, attracts, and on meeting annihilates. Same repels; opposite pairs. The reader may notice the shape of something here --- the rule that lets two things into one place only if they are opposite --- and I want to be careful to say that this is a shape and not yet a derivation.

\begin{figure}[htbp]
\centering
\begin{tikzpicture}[>=stealth]
  % ---- (a) same sense ----
  \begin{scope}[xshift=-3.1cm]
    \draw[dashed, black!45] (0,0) ellipse (2.15 and 1.15);
    \foreach \cx in {-0.85,0.85} {
      \draw[bbuBlue, line width=1.0pt] (\cx,0) circle (0.55);
      \draw[->, bbuBlue, line width=1.1pt] (\cx,0) ++(72:0.55) arc (72:108:0.55);
      \draw[->, bbuBlue, line width=1.1pt] (\cx,0) ++(252:0.55) arc (252:288:0.55);
      \fill[black!70] (\cx,0) circle (0.05);
    }
    \node[anchor=north, align=center, text width=4.6cm] at (0,-1.35)
      {\small (a) same sense: from outside, one circulation of twice the strength\\ $\oint = 2\kappa$ \quad energy $\propto 4$};
  \end{scope}
  % ---- (b) opposite sense ----
  \begin{scope}[xshift=3.1cm]
    \draw[dashed, black!45] (0,0) ellipse (2.15 and 1.15);
    \draw[bbuBlue, line width=1.0pt] (-0.85,0) circle (0.55);
    \draw[->, bbuBlue, line width=1.1pt] (-0.85,0) ++(72:0.55) arc (72:108:0.55);
    \draw[->, bbuBlue, line width=1.1pt] (-0.85,0) ++(252:0.55) arc (252:288:0.55);
    \fill[black!70] (-0.85,0) circle (0.05);
    \draw[bbuRed, line width=1.0pt] (0.85,0) circle (0.55);
    \draw[->, bbuRed, line width=1.1pt] (0.85,0) ++(108:0.55) arc (108:72:0.55);
    \draw[->, bbuRed, line width=1.1pt] (0.85,0) ++(288:0.55) arc (288:252:0.55);
    \fill[black!70] (0.85,0) circle (0.05);
    \node[anchor=north, align=center, text width=4.6cm] at (0,-1.35)
      {\small (b) opposite senses: from outside, nothing at all\\ $\oint = 0$ \quad energy $\propto 0$};
  \end{scope}
\end{tikzpicture}
\caption{Why two like circulations are a bad bargain. Energy goes as the square of circulation, and circulations add at a distance. (a)~Two of the same sense look, from far enough away, like one of twice the strength: four units of energy where the two of them apart cost two --- so a double quantum splits, as doubly-quantized vortices are observed to do. (b)~Two of opposite sense cancel at a distance, cost almost nothing, attract, and on meeting annihilate. Same and opposite are not symmetric. That asymmetry is the shape of the rule which admits two things to one place only in opposite senses --- the shape of it, and not yet the substance.}
\label{fig:circulationpair}
\end{figure}

There is a third fact, and it is the one that speaks to Chapter~7 rather than to Chapter~6. Rotate the helium harder and the vortices do not smear out and do not pile up: they arrange themselves into a lattice with a definite spacing, set by how hard you are rotating. Press harder still and the spacing shrinks, until the cores touch --- and then the arrangement gives way and something else happens. That is a strut being overwhelmed, in a bucket, on a bench. If exclusion is congestion rather than decree, then congestion can be \emph{beaten}, and a ladder of collapses is exactly what one should expect. The account I do not yet have would not only explain the strut; it would explain why there is a ladder at all.

And the same relation pays once more, in a place I had not thought to look for it. The third wound has this book's charges \emph{tumbling}: turning end over end, so that what a distant body meets is an average over orientations. Suppose that tumbling is not a device introduced to save Coulomb's law but the thing itself --- suppose it is what physics calls spin, taken literally.

Physics declines to take it literally, and has an arithmetic reason. An electron carries angular momentum of half of Planck's constant over $2\pi$; ask how fast a small ball would have to turn to carry that much, and the answers are ridiculous. At the old classical electron radius the surface must move at a hundred and seventy times the speed of light. At the size experiment now permits --- no bigger than the hundred-thousandth of a proton quoted in the fourth wound --- the figure is nearer fifty million times. So the textbooks conclude, reasonably, that spin is not rotation of anything: it is an intrinsic property, possessed and not explained. Another floor.

The reader can see where this book must go. \emph{Faster than light} is a violation only against our layer's light, and this chapter has twice been obliged to say that our $c$ is a local ordinance and the floor below runs faster --- once for the waves, and again for the correlations, where the fifth wound bought sub-layer channels outright. A circulation whose parts move at ten million times our light speed is unremarkable if it is a circulation of the layer below. The objection that removes literal spin from the standard picture simply does not bind here. That is not a proof that spin is rotation. It is the removal of the reason for saying it cannot be.

And the magnitude comes free, which surprised me. Take the confinement relation above --- a circulation of size $d$ carries momentum $\hbar/d$ --- and ask what angular momentum that is, about the centre: $L = p\,d = \hbar$. One unit. The $d$ cancels, and so does the mass: it does not matter how small the thing is or how heavy. A quantized circulation cannot help carrying angular momentum of order $\hbar$, because that is what a quantum of circulation \emph{is}. Spin, in this picture, is not a further property the electron happens to have. It is the stiffness of the last few pages, read around instead of across.

Two things it does not give, and I will not pretend otherwise. It gives one unit and not a half --- and a thing that merely rotates always carries whole units; the half is the tether's business rather than the turning's. And a rotating ball of charge has a magnetic moment in a fixed ratio to its spin, which for the electron is wrong by a factor of two. So there are now three twos owed: the two things a state will admit, the two turns that bring a tethered object home, and the two by which the electron's magnetism exceeds its turning. In the standard picture these are not three facts but one --- Dirac's equation delivers all of them together --- and I take that as the useful hint in the whole business. If they are one fact here too, then whoever supplies the packing rule will not be paying three debts. They will be paying one.


Three difficulties, stated plainly, because this is the least finished thing in the book. First, straight vortex lines are effectively two-dimensional objects when it comes to packing, and the five-thirds above wants \emph{compact} circulations --- closed loops, rings --- packed in three dimensions; the book's particles would have to be of that kind. Second, vortices can cross and swap ends. They are not perfectly impenetrable, and any account will have to say why matter's circulations are stricter than helium's. Third and hardest: a superfluid is a condensate, which is to say a crowd of things that \emph{like} to share a state, and getting the opposite behaviour out of structures living in it is not free. Nor have I derived the factor of two --- the rule permits two things in a state, of opposite sense, not one, and my figure shows a shape where a derivation should stand.

But the hardest of these --- and the half-unit of spin along with it, for they are the same difficulty wearing two hats --- has been done before, by people who were not trying to help this book. It is an established result that structures in a purely bosonic medium can be quantized as \emph{fermions}, carrying half-integer spin: Finkelstein and Rubinstein showed the possibility in 1968, and Skyrme's model of the nucleon --- a soliton in a field of bosons that comes out as a spin-one-half particle --- is the working example everyone knows. The mechanism has a homely demonstration: a thing tethered to its surroundings is not returned to itself by one full turn, because the tethers come back twisted, and needs a second turn to undo them. Something embedded in a medium is tethered by construction. So the last difficulty is not a wall. It is a place where somebody else has already been.

What I have, then, is not an account. It is a specification with two of its three parts already lying about in laboratory glassware: a mechanism that forces motion under confinement, in exactly the proportion the stiffness requires; and an arithmetic that makes like circulations separate and opposite ones pair. The missing part is the packing rule itself, with its factor of two --- and until somebody supplies it, the ladder of Chapter~7 stands on a strut whose nature is described rather than derived. I would rather print the specification than pretend to the account.

